%Sobreposição de oscilações
%Written by Tordar

\documentclass[12pt]{article}

\usepackage{graphicx}
\usepackage{pifont}
\usepackage{pstricks-add}
\usepackage{pst-infixplot}
\usepackage{pst-plot}

\pagestyle{empty} 

\begin{document}

\begin{pspicture}(3,-7)(18,7)
\psset{xunit=12mm,yunit=20mm}
%Plot the first oscilating term:
\rput{0}(0,3.5){
\infixtoRPN{sin(x*360/3.1415)*cos(180)}
\psplot[plotpoints=200,linecolor=red,linewidth=2pt]{0}{14}{\RPN}
\psline[linecolor=blue,linewidth=2pt]{->}(0,-1.5)(0,1.5)
\psline[linecolor=blue,linewidth=2pt]{->}(0,0)(14.5,0)
\rput{0}(15,0){\scalebox{1.5}{$t$}}
}
%Plot the first two oscilating terms:
\infixtoRPN{sin(x*360/3.1415)*cos(180)-sin(x*720/3.1415)/2}
\psplot[plotpoints=200,linecolor=magenta,linewidth=2pt]{0}{14}{\RPN}
\psline[linecolor=blue,linewidth=2pt]{->}(0,-1.5)(0,1.5)
\psline[linecolor=blue,linewidth=2pt]{->}(0,0)(14.5,0)
\rput{0}(15,0){\scalebox{1.5}{$t$}}
%Plot the first three oscilating terms:
\rput{0}(0,-3.5){
\infixtoRPN{sin(x*360/3.1415)*cos(180)-sin(x*720/3.1415)/2-sin(x*1080/3.1415)/3}
\psplot[plotpoints=200,linecolor=green,linewidth=2pt]{0}{14}{\RPN}
\psline[linecolor=blue,linewidth=2pt]{->}(0,-1.5)(0,1.5)
\psline[linecolor=blue,linewidth=2pt]{->}(0,0)(14.5,0)
\rput{0}(15,0){\scalebox{1.5}{$t$}}
}
\rput{0}(0,-7){
\infixtoRPN{sin(x*360/3.1415)*cos(180)-sin(x*720/3.1415)/2-sin(x*1080/3.1415)/3-sin(x*1440/3.1415)/4}
\psplot[plotpoints=200,linecolor=green,linewidth=2pt]{0}{14}{\RPN}
\psline[linecolor=blue,linewidth=2pt]{->}(0,-1.5)(0,1.5)
\psline[linecolor=blue,linewidth=2pt]{->}(0,0)(14.5,0)
\rput{0}(15,0){\scalebox{1.5}{$t$}}
}
\end{pspicture}

\end{document}
